% ==============================================================================
% Figura: Diagrama de Bloques Atómico — Estado FURUTA_STATE_MOTOR_PID_VEL (1 kHz)
% Estrategia: 1 función = 1 bloque atómico, solo nombre corto en cajas, solo
% símbolos puros de señal en cables (sin ganancias ni ecuaciones en bloques).
% Fiel a furuta_control.c (líneas 362-369), pid_controller.c, l298n_motor.c
% y as5600.c.
% Archivo modular listo para \input{} en IEEEtran dos columnas (figure* [!t]).
% Bibliotecas TikZ requeridas: arrows.meta, positioning, fit, calc,
%                            shapes.geometric, backgrounds, babel, bending.
% ==============================================================================
\begin{figure*}[!t]
\centering
% --- Paleta semántica armónica (idéntica a fig:diagrama_bloques_control) ---
\definecolor{blkInk}{HTML}{1E293B}
\definecolor{blkPendEdge}{HTML}{1B4965}
\definecolor{blkPendFill}{HTML}{EAF4FB}
\definecolor{blkArmEdge}{HTML}{0F766E}
\definecolor{blkArmFill}{HTML}{E6F7F5}
\definecolor{blkCondEdge}{HTML}{B45309}
\definecolor{blkCondFill}{HTML}{FEF6EB}
\definecolor{blkDistEdge}{HTML}{7E2268}
\definecolor{blkDistFill}{HTML}{FBF0F8}
\definecolor{blkActEdge}{HTML}{A63429}
\definecolor{blkActFill}{HTML}{FDF0EE}
\definecolor{blkPlantEdge}{HTML}{243447}
\definecolor{blkPlantFill}{HTML}{ECEFF4}
\definecolor{blkZoneZ1}{HTML}{33658A}
\definecolor{blkZoneL2}{HTML}{A63429}
\definecolor{blkZoneL3}{HTML}{2A6B43}

% --- Parámetros centralizados de ritmo geométrico (Layout by Construction) ---
\def\bvBoxH{0.86cm}
\def\bvSubH{0.80cm}
\def\bvSumSz{0.58cm}
\def\bvHopRad{0.095cm}
\def\bvRefLeadX{1.48cm}
\def\bvOutLeadX{1.34cm}
\def\bvErrSplitX{0.64cm}
\def\bvErrTapX{0.76cm}
\def\bvSumCGapX{1.02cm}
\def\bvSatGapX{0.54cm}
\def\bvBranchGapY{0.56cm}
\def\bvPropGapY{0.46cm}
\def\bvZoneGapX{1.62cm}
\def\bvActGapX{0.54cm}
\def\bvVmGapX{0.78cm}
\def\bvSensGapX{0.44cm}
\def\bvZoneThreeDropY{0.68cm}
\def\bvEncBackX{0.58cm}

\resizebox{0.97\textwidth}{!}{%
\begin{tikzpicture}[
    >=Stealth,
    font=\sffamily,
    % --- Estilos de bloques atómicos (1 función = 1 bloque, solo nombre corto) ---
    baseblk/.style={
        rectangle,
        rounded corners=3.5pt,
        align=center,
        inner xsep=3.5pt,
        inner ysep=3pt,
        text=blkInk,
        font=\normalsize\sffamily\bfseries,
        minimum height=\bvBoxH
    },
    armblk/.style={baseblk, draw=blkArmEdge, line width=1.1pt, fill=blkArmFill, text width=2.16cm},
    condblk/.style={baseblk, draw=blkCondEdge, line width=1.1pt, fill=blkCondFill, text width=1.52cm},
    distblk/.style={baseblk, draw=blkDistEdge, line width=1.05pt, fill=blkDistFill, text width=2.16cm, minimum height=\bvSubH},
    actblk/.style={baseblk, draw=blkActEdge, line width=1.1pt, fill=blkActFill, text width=1.74cm},
    plantblk/.style={baseblk, draw=blkPlantEdge, line width=1.35pt, fill=blkPlantFill, text width=1.46cm},
    sensTblk/.style={baseblk, draw=blkArmEdge, line width=1.1pt, fill=blkArmFill, text width=2.24cm},
    sum/.style={
        circle,
        draw=blkInk,
        line width=1.05pt,
        fill=white,
        minimum size=\bvSumSz,
        inner sep=0pt,
        font=\small\bfseries,
        text=blkInk
    },
    dot/.style={circle, fill=#1, inner sep=1.55pt},
    % --- Flechas y trazos canónicos ---
    wire/.style={line width=1.0pt, draw=blkInk},
    sig/.style={-{Stealth[length=4.8pt, width=3.5pt]}, line width=1.0pt, draw=blkInk},
    sigshort/.style={-{Stealth[length=3.9pt, width=2.9pt]}, line width=1.0pt, draw=blkInk},
    fbT/.style={-{Stealth[length=4.4pt, width=3.2pt]}, line width=0.95pt, draw=blkArmEdge},
    lineT/.style={line width=0.95pt, draw=blkArmEdge},
    adapt/.style={-{Stealth[length=3.9pt, width=2.9pt]}, line width=0.95pt, densely dashed, draw=blkDistEdge},
    lineDist/.style={line width=0.95pt, densely dashed, draw=blkDistEdge},
    % --- Contenedores y etiquetas ---
    zonebox/.style={
        rectangle,
        rounded corners=5pt,
        draw=#1,
        densely dashed,
        line width=0.95pt,
        fill=#1,
        fill opacity=0.035,
        text opacity=1
    },
    zonetitle/.style={
        rectangle,
        rounded corners=3pt,
        draw=#1,
        line width=0.8pt,
        fill=white,
        inner xsep=5.5pt,
        inner ysep=2.2pt,
        font=\small\sffamily\bfseries,
        text=#1
    },
    signlbl/.style={font=\footnotesize\bfseries, text=blkInk, inner sep=1pt},
    wirelbl/.style={font=\small\sffamily, text=blkInk, inner sep=1.5pt}
]

% ==============================================================================
% ZONA 1: CONTROL PID DE VELOCIDAD (PARALELO) — ÚNICO ANCLAJE (0,0)
% ==============================================================================
\node[sum] (sig_v) at (0,0) {$\Sigma_v$};
\coordinate[right=\bvErrSplitX of sig_v.east] (split_e);

\node[armblk, right=\bvErrTapX of split_e] (int_tustin) {Integrador\\Tustin};
\node[distblk, above=\bvBranchGapY of int_tustin] (anti_windup) {Anti-Windup};
\node[armblk, above=\bvPropGapY of anti_windup] (prop_blk) {Proporcional};
\node[armblk, below=\bvBranchGapY of int_tustin] (deriv_blk) {Derivativo\\Filtrado};

\node[sum, right=\bvSumCGapX of int_tustin] (sig_c) {$\Sigma_c$};
\node[condblk, right=\bvSatGapX of sig_c] (sat_vel) {Sat. Vel.};

\coordinate (z1_top_pad)   at ([yshift=0.46cm]prop_blk.north);
\coordinate (z1_bot_pad)   at ([yshift=-0.28cm]deriv_blk.south);
\coordinate (z1_left_pad)  at ([xshift=-0.56cm]sig_v.west);
\coordinate (z1_right_pad) at ([xshift=0.28cm]sat_vel.east);

\begin{scope}[on background layer]
    \node[zonebox=blkZoneZ1, inner sep=0pt,
          fit=(z1_top_pad)(z1_bot_pad)(z1_left_pad)(z1_right_pad)] (zone1) {};
\end{scope}
\node[zonetitle=blkZoneZ1, anchor=west] at ([xshift=0.35cm]zone1.north west)
    {Control PID de Velocidad (Paralelo)};

% ==============================================================================
% ZONA 2: ACTUACIÓN Y MOTOR (CASCADA)
% ==============================================================================
\node[actblk, right=\bvZoneGapX of sat_vel] (banda) {Banda\\Silenciosa};
\node[actblk, right=\bvActGapX of banda] (rampa) {Rampa\\L298N};
\node[actblk, right=\bvActGapX of rampa] (pwm) {Modulador\\PWM};
\node[actblk, right=\bvActGapX of pwm] (puente) {Puente H\\L298N};
\node[plantblk, right=\bvVmGapX of puente] (motor) {Motor\\80:1};

\coordinate (l2_top_pad)   at ([yshift=0.42cm]banda.north);
\coordinate (l2_bot_pad)   at ([yshift=-0.28cm]banda.south);
\coordinate (l2_left_pad)  at ([xshift=-0.28cm]banda.west);
\coordinate (l2_right_pad) at ([xshift=0.28cm]motor.east);

\begin{scope}[on background layer]
    \node[zonebox=blkZoneL2, inner sep=0pt,
          fit=(l2_top_pad)(l2_bot_pad)(l2_left_pad)(l2_right_pad)] (zoneL2) {};
\end{scope}
\node[zonetitle=blkZoneL2, anchor=west] at ([xshift=0.32cm]zoneL2.north west)
    {Actuaci\'on y Motor (Cascada)};

% ==============================================================================
% ZONA 3: SENSADO BASE Y LPF IIR 1 (DERECHA -> IZQUIERDA)
% ==============================================================================
\coordinate (l3_y) at ([yshift=-\bvZoneThreeDropY]zone1.south);
\coordinate (as5600_east_x) at ([xshift=-\bvEncBackX]motor.south);

\node[sensTblk, anchor=east] (as5600_base) at (as5600_east_x |- l3_y) {AS5600\\Base};
\node[sensTblk, left=\bvSensGapX of as5600_base] (desenv) {Desenvoltura};
\node[sensTblk, left=\bvSensGapX of desenv] (deriv_t) {Derivador $\theta$};
\node[sensTblk, left=\bvSensGapX of deriv_t] (lpf1_t) {LPF IIR 1};

\coordinate (l3_top_pad)   at ([yshift=0.42cm]lpf1_t.north);
\coordinate (l3_bot_pad)   at ([yshift=-0.30cm]lpf1_t.south);
\coordinate (l3_left_pad)  at ([xshift=-0.28cm]lpf1_t.west);
\coordinate (l3_right_pad) at ([xshift=0.28cm]as5600_base.east);

\begin{scope}[on background layer]
    \node[zonebox=blkZoneL3, inner sep=0pt,
          fit=(l3_top_pad)(l3_bot_pad)(l3_left_pad)(l3_right_pad)] (zoneL3) {};
\end{scope}
\node[zonetitle=blkZoneL3, anchor=west] at ([xshift=0.32cm]zoneL3.north west)
    {Sensado Base y LPF IIR 1};

% ==============================================================================
% CONEXIONES ZONA 1 (PID PARALELO + SUPERVISIÓN ANTI-WINDUP)
% ==============================================================================
\coordinate[left=\bvRefLeadX of sig_v.west] (ref_in);

% Entrada de referencia omega_ref (centrada fuera del contenedor Zona 1)
\draw[sig] (ref_in) -- (sig_v.west);
\node[wirelbl, anchor=south]
    at ($(ref_in)!0.50!(zone1.west |- sig_v.west) + (0, 2pt)$) {$\omega_{\mathrm{ref}}$};
\node[signlbl, anchor=south] at ([xshift=-0.14cm, yshift=0.07cm]sig_v.west) {$+$};

% Error e_omega y bifurcación paralela hacia P, I y Anti-Windup
\draw[wire] (sig_v.east) -- (split_e)
    node[wirelbl, midway, above=2pt] {$e_\omega$};
\node[dot=blkInk] at (split_e) {};

\draw[sigshort] (split_e) -- (int_tustin.west);

\coordinate (aw_tap) at (split_e |- anti_windup.west);
\draw[wire] (split_e) -- (aw_tap);
\node[dot=blkInk] at (aw_tap) {};

\draw[sigshort] (aw_tap) -- (anti_windup.west)
    node[wirelbl, midway, above=2pt] {$e_\omega$};
\draw[sig, rounded corners=4pt] (aw_tap) |- (prop_blk.west);

% Supervisión Anti-Windup -> clamp hacia Integrador Tustin
\draw[adapt] (anti_windup.south) -- (int_tustin.north)
    node[wirelbl, text=blkDistEdge, midway, right=3pt] {$\mathrm{clamp}$};

% Convergencia paralela en Sigma_c
\draw[sig, rounded corners=4pt] (prop_blk.east) -| (sig_c.north)
    node[wirelbl, pos=0.25, above=2pt] {$u_p$};
\node[signlbl, anchor=south west] at ([xshift=0.08cm, yshift=0.07cm]sig_c.north) {$+$};

\draw[sigshort] (int_tustin.east) -- (sig_c.west)
    node[wirelbl, pos=0.38, above=2pt] {$u_i$};
\node[signlbl, anchor=south] at ([xshift=-0.14cm, yshift=0.07cm]sig_c.west) {$+$};

\draw[sig, rounded corners=4pt] (deriv_blk.east) -| (sig_c.south)
    node[wirelbl, pos=0.25, below=2pt] {$u_d$};
\node[signlbl, anchor=north west] at ([xshift=0.08cm, yshift=-0.07cm]sig_c.south) {$+$};

% Salida de Sigma_c -> Sat. Vel. y realimentación sat -> Anti-Windup (con puente sobre u_p)
\draw[sigshort] (sig_c.east) -- (sat_vel.west);

\coordinate (sat_hop_x) at (sig_c.center);
\coordinate (sat_bus_y) at (anti_windup.east);
\draw[lineDist, rounded corners=4pt]
    (sat_vel.north) |- ($(sat_hop_x |- sat_bus_y)+(\bvHopRad,0)$)
    node[wirelbl, text=blkDistEdge, pos=0.72, above=2pt] {$\mathrm{sat}$};
\draw[lineDist]
    ($(sat_hop_x |- sat_bus_y)+(\bvHopRad,0)$)
    arc[start angle=0, end angle=180, radius=\bvHopRad];
\draw[adapt]
    ($(sat_hop_x |- sat_bus_y)+(-\bvHopRad,0)$) -- (anti_windup.east);

% ==============================================================================
% CONEXIONES ZONA 2 (CASCADA L298N -> MOTOR 80:1 -> SALIDA omega_m)
% ==============================================================================
\draw[sig] (sat_vel.east) -- (banda.west);
\node[wirelbl, anchor=south]
    at ($(zone1.east |- sat_vel.east)!0.50!(zoneL2.west |- banda.west) + (0, 2pt)$) {$u_v[k]$};

\draw[sigshort] (banda.east) -- (rampa.west);
\draw[sigshort] (rampa.east) -- (pwm.west);
\draw[sigshort] (pwm.east) -- (puente.west);
\draw[sigshort] (puente.east) -- (motor.west)
    node[wirelbl, midway, above=2pt] {$V_m$};

\coordinate[right=\bvOutLeadX of motor.east] (omega_out);
\draw[sig] (motor.east) -- (omega_out);
\node[wirelbl, anchor=south]
    at ($(zoneL2.east |- motor.east)!0.50!(omega_out) + (0, 2pt)$) {$\omega_m$};

% ==============================================================================
% CONEXIONES ZONA 3 (ENCODER BASE -> LPF IIR 1 -> RETORNO \dot{\theta}_{f1})
% ==============================================================================
\draw[fbT, rounded corners=4pt] (motor.south) |- (as5600_base.east);
\node[wirelbl, text=blkArmEdge, anchor=west]
    at ([xshift=0.08cm]$(motor.south |- zoneL2.south)!0.48!(motor.south |- zoneL3.north)$) {$\theta(t)$};

\draw[fbT] (as5600_base.west) -- (desenv.east);
\draw[fbT] (desenv.west) -- (deriv_t.east);
\draw[fbT] (deriv_t.west) -- (lpf1_t.east);

% Retorno limpio de \dot{\theta}_{f1} por debajo de Zona 1 (CERO cruces)
\coordinate (vel_tap) at (sig_v.south |- deriv_blk.west);

\draw[lineT, rounded corners=4pt]
    (lpf1_t.west) -| (vel_tap);
\node[wirelbl, text=blkArmEdge, anchor=north]
    at ($(lpf1_t.west)!0.42!(sig_v.south |- lpf1_t.west) + (0, -0.04cm)$) {$\dot{\theta}_{f1}$};

\node[dot=blkArmEdge] at (vel_tap) {};

\draw[fbT] (vel_tap) -- (sig_v.south);
\node[signlbl, text=blkArmEdge, anchor=north west]
    at ([xshift=0.08cm, yshift=-0.07cm]sig_v.south) {$-$};
\node[wirelbl, text=blkArmEdge, anchor=east]
    at ([xshift=-0.08cm]$(vel_tap)!0.50!(sig_v.south)$) {$\dot{\theta}_{f1}$};

\draw[fbT] (vel_tap) -- (deriv_blk.west)
    node[wirelbl, text=blkArmEdge, pos=0.56, above=2pt] {$-\dot{\theta}_{f1}$};

\end{tikzpicture}%
}
\caption{Diagrama de bloques at\'omico del estado \texttt{FURUTA\_STATE\_MOTOR\_PID\_VEL} a \qty{1.0}{\kilo\hertz} (\texttt{PID\_Compute} en \texttt{pid\_controller.c} y \texttt{L298N\_SetVoltage} en \texttt{l298n\_motor.c}). Muestra el esquema paralelo discretizado por Tustin con derivada directa sobre la medici\'on negada y supervisi\'on \textit{anti-windup} por congelamiento condicional (\texttt{clamp}) ante saturaci\'on en $\pm\qty{12.0}{\volt}$, la cascada de potencia L298N hacia el motor $80:1$ y la rama \'unica de sensado base AS5600 hasta \texttt{LPF IIR 1}.}
\label{fig:bloques_velocidad}
\end{figure*}
